\scanpage{190}
\mspage{139}
par $X\mapsto X^0$ et par $X\mapsto\Sigma(X)$, et en tenant compte des hom fonctoriels en $X$
\diagram{\begin{tikzcd}[ampersand replacement=\&,row sep=tiny,column sep=small]
\&\Sigma(X)\arrow[dl]\arrow[dr]\&\\X\&\&X^0\rlap{,}
\end{tikzcd}}
à trois diagrammes du type $\ast\ast)$, i.e. modelés sur la catégorie opposée de
\diagram{\begin{tikzcd}[ampersand replacement=\&,row sep=small,column sep=small]
\&\&0\&\\
\Gamma={}\&1\arrow[ru,"\alpha_1"]\&\&2\arrow[lu,"\alpha_2"']\\
\&\&3\arrow[lu,"i_1"']\arrow[ru,"i_2"]\&\\
\&\&4\arrow[u,"\unsure{\varphi|\rho}"]\&
\end{tikzcd}}
% slanted
\margin{On \deleted{\unsure{impose}} a $\alpha_1i_1\ne\alpha_2i_2$, mais la relation $\alpha_1i_1\varphi=\alpha_2i_{\unsure{2|1}}\varphi$}
(i.e. \tikz[baseline=(a.base)]{\node[draw,ellipse,inner sep=2pt] (a) {$\Delta^1\times\Delta^1$};} auquel on a rajouté un objet initial),
% slanted
\margin{faux, car le carré n'est pas commutatif}
savoir
\[
\widehat Q,\qquad\widehat{\Sigma(Q)},\qquad\widehat{Q^0}
\]
et deux homomorphismes de diagrammes ($\in\underline{\mathrm{Hom}}(\Gamma,\mathit{Cat})$)
\diagram{\begin{tikzcd}[ampersand replacement=\&,column sep=small,row sep=small]
\&\&\widehat{\Sigma(Q)}\arrow[dl]\arrow[dr]\&\\
(\ast\ast\ast)\&\widehat Q\&\&\widehat{Q^0}
\end{tikzcd}}
(déduit de
\diagram{\begin{tikzcd}[ampersand replacement=\&,column sep=small,row sep=tiny]
\&\Sigma(Q)\arrow[dl]\arrow[dr]\&\\Q\&\&Q^0
\end{tikzcd}}
), lesquels sont des $W$ terme à terme (i.e. \deleted{dans} \unsure{dans} $W_\Gamma\subset\mathrm{Fl}(\underline{\mathrm{Hom}}(\Gamma,\mathit{Cat}))$). Explicitement, en les sommets $i=0,1,2$ \inserted{resp. $4$} de $\Gamma$, ceci n'est autre que
\diagram{\begin{tikzcd}[ampersand replacement=\&,column sep=tiny,row sep=tiny]
\&\&\Sigma(X_i)\arrow[dl,"p_{X_i}"']\arrow[dr,"q_{X_i}"]\&\&\&\&\Sigma Y\arrow[dl,"p_Y"']\arrow[dr,"q_Y"]\&\\
D_i=\Biggl(\&X_i\&\&X_i^0\Biggr)\&\text{resp.}\&Y\&\&Y^0
\end{tikzcd}}
en le sommet $3$ c'est
\diagram{\begin{tikzcd}[ampersand replacement=\&,column sep=small,row sep=small]
\&\int\Sigma X_\ast\arrow[dl,"\int p_{X_\ast}"']\arrow[dr,"\int q_{X_\ast}"]\&\\
\mathscr X=\int X_\ast\&\&\int X_\ast^0\simeq\widetilde{\mathscr X}^0
\end{tikzcd}}
obtenu par \textquotedblleft intégration\textquotedblright\ des \deleted{\unsure{flèch.}} flèches
\diagram{\begin{tikzcd}[ampersand replacement=\&,column sep=tiny,row sep=tiny]
\&\deleted{\unsure{\Sigma X_\ast}}\arrow[dl,"\deleted{\unsure{p_{X_\ast}}}"',"\backslash"{pos=.50,anchor=center}]\\
\deleted{\unsure{X_\ast}}\&
\end{tikzcd}}
\diagram{\begin{tikzcd}[ampersand replacement=\&,column sep=small,row sep=tiny]
\&\Sigma X_\ast\arrow[dl,"p_{X_\ast}"']\arrow[dr,"q_{X_\ast}"]\&\\
X_\ast\&\&X_\ast^0 .
\end{tikzcd}}
% archive sheet number: 190
